Using this knowledge we can derive an ideal low pass filter.
\[z_i = \begin{cases} 0, \;i \in (k;n-k) \\ 1, \; \text{otherwise} \end{cases}\]

This is an impulse reponse of the filter cutting all the frequencies above the k-th one. Recall that the $N-k$ element in DFT vector corresponds to the negative oscillation with the frequency $k$, due to cyclicity of $e^{\frac{2\pi i j x}{N}}$.

Using the formula for \ref{eq:IDFT}:
\[\begin{aligned}
		y_{j} = \sum_{l=0}^{N-1} F_{jl} z_k
		= \sum_{l=0}^{k} F_{jl} + \sum_{l=N-k}^{N-1} F_{jl}
		= \frac{1}{\sqrt{N}} \left(\sum_{l=0}^{k} \omega_N^{jl} + \sum_{l=N-k}^{N-1} \omega_N^{jl} \right)
	\end{aligned}\]

Substitute $l = N-m-1$ in the second sum:
\[\begin{aligned}
		\frac{1}{\sqrt{N}} \left(\sum_{l=0}^{k} \omega_N^{jl} + \sum_{l=N-k}^{N-1} \omega_N^{jl} \right) =
		\frac{1}{\sqrt{N}} \left(\sum_{l=0}^{k} \omega_N^{jl} + \sum_{m=0}^{k-1} \omega_N^{j(N-m-1)} \right)
		 & =                                               \\ \frac{1}{\sqrt{N}} \left(\sum_{l=0}^{k} \omega_N^{jl} + \sum_{m=0}^{k-1} \omega_N^{j(N-m-1)} \right)
		=\frac{1}{\sqrt{N}} \left(\sum_{l=0}^{k} \omega_N^{jl} + \sum_{m=0}^{k-1} \omega_N^{jN} \omega_N^{-j} \omega_N^{-jm} \right)
		 & \underset{\ref{lm:roots of unity cyclicity}}{=} \\ \frac{1}{\sqrt{N}} \left(\omega_N^{jk} + \sum_{l=0}^{k-1} \omega_N^{jl} + \omega_N^{-j} \sum_{m=0}^{k-1} \omega_N^{-jm} \right)
	\end{aligned}\]

Let $j\neq 0$. Using the sum of geometric series:
\[\sum_{l=0}^{k-1} \omega_N^{jl} = \sum_{l=0}^{k-1} {\left(\omega_N^{j}\right)}^l = \frac{1 - {\left(\omega_N^{j}\right)}^k}{1 - \omega_N^j}\]

Similarly:
\[\begin{aligned}
		\omega_N^{-j}\sum_{m=0}^{k-1} \omega_N^{-jm}
		= \omega_N^{-j}\sum_{m=0}^{k-1} {\left(\omega_N^{-j}\right)}^m
		= \frac{\omega_N^{-j} - \omega_N^{-j}{\left(\omega_N^{-j}\right)}^k}{1 - \omega_N^{-j}}
		 & = \\ \frac{\omega_N^{-j} - \omega_N^{-j}{\left(\omega_N^{-j}\right)}^k}{1 - \omega_N^{-j}} \frac{\omega_N^j}{\omega_N^j}
		= \frac{1 - {\left(\omega_N^{-j}\right)}^k}{\omega_N^{j} - 1}
		= \frac{{\left(\omega_N^{-j}\right)}^k - 1}{1 - \omega_N^{j}}
	\end{aligned}\]

By combining the terms together:
\[
	\sum_{l=0}^{k-1} \omega_N^{jl} + \omega_N^{-j} \sum_{l=0}^{k-1} \omega_N^{-jl}
	= \frac{1 - {\left(\omega_N^{j}\right)}^k}{1 - \omega_N^j} + \frac{{\left(\omega_N^{-j}\right)}^k - 1}{1 - \omega_N^{j}}
	= \frac{\omega^{-jk} - \omega^{jk}}{1 - \omega_N^{j}}\]

This yields the following expression for $y_j$, when $j\neq 0$:
\[y_{j}  = \frac{1}{\sqrt{N}} \left(\omega_N^{jk} + \frac{\omega_N^{-jk} - \omega_N^{jk}}{1 - \omega_N^{j}} \right)\]

If $j = 0$, then all the summands are equal to one, so get just $(2k+1)/\sqrt{N}$.
By the Theorem $\ref{thm:discrete convolution theorem}$:
\[\frac{1}{\sqrt N}(F(x\circledast y)) = Fx \odot Fy\]
which allows us to implement the low pass filter using its impulse response. That is the point, where theory meets the real world example and makes use of the abstract theorem. In general this concept can be generalaized to any signal, which can be used, for instance for reverb implementation if we have an impulse response of a room (a recording of echo induced by a short impact).

Now we take a closer look on the expression we have derived above:
\[
	e^{-ix}-e^{ix} = 2\imag(e^{-ix}) = 2i\sin(-x) = -2i\sin(x)\]
\[
	1-e^{ix}
	= e^{\frac{ix}{2}}\left(e^{-\frac{ix}{2}} - e^{\frac{ix}{2}}\right)
	= -2i\cdot e^{\frac{ix}{2}} \sin \left(\frac{x}{2}\right)\]

Therefore:
\[
	1 - \omega_N^{j} = 1-e^{2\pi i\frac{j}{N}}
	= -2i\cdot e^{\frac{\pi i\cdot j}{N}} \sin \left(\frac{\pi j}{N}\right)\]
\[
	\omega_N^{-jk} - \omega_N^{jk}
	= e^{-2\pi i\frac{kj}{N}}-e^{2\pi i\frac{kj}{N}}
	= -2i \sin\left(\frac{2\pi kj}{N}\right)\]
\[
	\frac{\omega_N^{-jk} - \omega_N^{jk}}{1 - \omega_N^{j}}
	= \frac{-2i \sin\left(\frac{2\pi kj}{N}\right)}{-2i\cdot e^{\frac{\pi i\cdot j}{N}} \sin \left(\frac{\pi j}{N}\right)}
	= \frac{e^{-\frac{\pi i\cdot j}{N}} \sin\left(\frac{2\pi kj}{N}\right)}{\sin \left(\frac{\pi j}{N}\right)}\]

With $x_j = 2\pi j / N$ we have:
\[
	\frac{e^{-\frac{\pi i\cdot j}{N}} \sin\left(\frac{2\pi kj}{N}\right)}{\sin \left(\frac{\pi j}{N}\right)}
	= \frac{e^{-ix_j/2} \sin\left(kx_j\right)}{\sin \left(x_j/2\right)}
	= \frac{\cos(x_j/2) \sin(kx_j)}{\sin(x_j/2)} - i \sin\left(kx_j\right)\]

Add the term $\omega_N^{jk}$ we had previously:
\begin{equation}
	\begin{aligned}
		\omega_N^{jk} + \frac{\cos(x_j/2) \sin(kx_j)}{\sin(x_j/2)} - i \sin(kx_j)
		 & = \\ \cos(kx_j) + i \sin(kx_j) + \frac{\cos(x_j/2) \sin(kx_j)}{\sin(x_j/2)} - i \sin(kx_j)
		 & = \\ \cos(kx_j) + \frac{\cos(x_j/2) \sin(kx_j)}{\sin(x_j/2)}
		 & = \\ \frac{\cos(kx_j) \sin(x_j/2) + \cos(x_j/2) \sin(kx_j)}{\sin(x_j/2)}
	\end{aligned}
	\label{eq:lpf before trig id}
\end{equation}

Using the equality $\sin(x+y) = \sin(x)\cos(y) + \sin(y)\cos(x)$:
\[\eqref{eq:lpf before trig id} =\frac{\sin((k+1/2)x_j)}{\sin(x_j/2)}\]
\[y_j = \begin{cases}
		\frac{1}{\sqrt{N}}\frac{\sin((k+1/2)x_j)}{\sin(x_j/2)}, \;  j\in \{1, \dots, N-1\} \\
		(2k+1)/\sqrt{N}, \; j = 0
	\end{cases}\]

Miraculously we got a continuous Dirichlet kernel, discretized at points $x_j$. One can show, that the impulse response of Dirichlet kernel in continuous space acts in the exact same way, namely cuts the frequencies above the k-th one. As we had seen in equation \eqref{eq:IDFT unnormalized}, $z_k = \sum_{n \in \Z} \hat{f}(k+nN)$, so given that $\widehat{D_n}(k) = \begin{cases}
		1, k \leq n \\
		0, \; \text{otherwise}
	\end{cases}$, this result was not completely unexpected.